\documentclass[11pt]{article}

\usepackage[a4paper, margin=1in]{geometry}
\usepackage{amsmath}
\usepackage{amssymb}
\usepackage{graphicx}
\usepackage{xcolor}
\usepackage{hyperref}
\usepackage{enumitem}

\hypersetup{colorlinks=true, linkcolor=black, urlcolor=blue}

\title{}
\date{}

\begin{document}

% ================= TITLE PAGE =================
\begin{center}
{\LARGE \bfseries Designing a Rational and Intelligent PACS System}\\[0.4em]
{\large Applying Agent-Based AI Concepts to Medical Image Management}
\end{center}

\vspace{0.5cm}

\begin{center}
\begin{tabular}{ll}
\textbf{Name:} & Devesh \\
\textbf{Roll Number:} & 24111014 \\
\textbf{Department:} & Biomedical Engineering \\
\textbf{Institution:} & NIT Raipur \\
\textbf{Course:} & AI \& ML \\
\textbf{Faculty:} & Prof. Saurabh Gupta \\
\textbf{Date:} & \today \\
\end{tabular}
\end{center}

\vspace{1cm}
\hrule
\vspace{0.5cm}

% ================= INTRODUCTION =================
\section*{Introduction}

A PACS (Picture Archiving and Communication System) is basically the system hospitals use to store, retrieve, and manage medical images like X-rays, CT scans, MRIs, and ultrasounds. This assignment looks at PACS through the lens of the agent-based AI concepts we've been covering -- rationality, intelligence, semantics, environment, complexity, and agent behaviour -- and tries to work out what a "smart" version of this system would actually look like.

% ================= Q1 =================
\section{How can we make a rational PACS system?}

A rational agent is one that picks the action expected to give the best outcome, based on the information it has. For a PACS system, that would mean:

\begin{itemize}[noitemsep]
    \item Store and organize medical images like X-rays, CT scans, MRI and ultrasound.
    \item Use patient history and image data to make decisions based on available information.
    \item Use AI models to detect abnormalities and support doctors in diagnosis.
    \item Choose the action that gives the best expected result, such as flagging a suspicious scan.
    \item Keep patient safety and accuracy as the main priorities.
\end{itemize}

% ================= Q2 =================
\section{How can we make it an intelligent PACS system so that it can pass the Turing Test?}

For the system to feel "intelligent" rather than just automated, it would need to behave more like a human expert interacting with the data:

\begin{itemize}[noitemsep]
    \item Add AI-based image analysis for detecting diseases or abnormalities.
    \item Allow the system to understand simple natural-language queries from doctors.
    \item Make it explain its findings instead of only giving a prediction.
    \item Learn from previous cases and improve its performance.
    \item Provide responses similar to how a medical professional would interpret and discuss an image.
    \item Combine medical images with patient information for better decisions.
\end{itemize}

% ================= Q3 =================
\section{What kind of semantics should we have in our PACS system?}

Semantics here basically means giving the raw data actual meaning the system can reason about:

\begin{itemize}[noitemsep]
    \item Medical terminology such as organs, diseases, lesions and anatomical structures.
    \item Image-related semantics such as modality, body part, image quality and abnormality.
    \item Patient-related information such as age, symptoms and clinical history.
    \item Relationships between findings, for example, "lung nodule $\rightarrow$ possible tumor."
    \item Standard medical coding and terminology so that different hospitals and systems can understand the data.
\end{itemize}

% ================= Q4 =================
\section{What defines the environment of the PACS system?}

\begin{itemize}[noitemsep]
    \item Medical images from X-ray, CT, MRI, ultrasound and other imaging devices.
    \item Patient medical records and clinical history.
    \item Radiologists, doctors, technicians and other healthcare professionals.
    \item Hospital information systems and electronic health records.
    \item Image acquisition devices and medical databases.
    \item The environment is partly observable because the system may not have all the patient's information.
    \item It is dynamic because new images, reports and patient information can be added over time.
\end{itemize}

% ================= Q5 =================
\section{How can we analyze the complexity of the PACS system?}

\begin{itemize}[noitemsep]
    \item Number and size of medical images being processed.
    \item Number of patients and records stored in the system.
    \item Complexity of AI models used for image analysis.
    \item Processing time required for diagnosis or image retrieval.
    \item Memory and storage requirements.
    \item Number of users and medical devices connected to the system.
    \item Complexity also increases when the system has to combine images, patient history and clinical information.
\end{itemize}

% ================= Q6 =================
\section{In which situations will our PACS system work as a deterministic agent and stochastic agent?}

\textbf{Deterministic agent:} when the same input always produces the same output, such as retrieving a particular patient's MRI using a patient ID.

\textbf{Stochastic agent:} when the result involves uncertainty, such as predicting whether an X-ray contains pneumonia.

AI-based diagnosis is generally stochastic because medical images can be unclear and different diseases can have similar appearances. So image classification and disease prediction can be treated as probabilistic decision-making rather than a fixed lookup.

% ================= Q7 =================
\section{What kind of tasks should we give to our PACS system?}

\begin{itemize}[noitemsep]
    \item Automatically classify and organize medical images.
    \item Detect abnormalities in X-rays, CT scans and MRI images.
    \item Highlight suspicious regions for the radiologist.
    \item Compare current images with previous patient scans.
    \item Retrieve relevant patient images quickly.
    \item Generate preliminary reports from detected findings.
    \item Prioritize urgent or abnormal cases.
    \item Check image quality before analysis.
    \item Assist doctors in making decisions, rather than completely replacing them.
\end{itemize}


\end{document}
